\documentclass{article}
 
\usepackage{booktabs}
\usepackage{tabularx}
\usepackage{array}
\usepackage{caption}
\usepackage[a4paper, margin=2cm]{geometry}
 
\begin{document}
\section{Conjunto de Dados}

O conjunto de dados utilizado neste estudo é composto por informações referentes aos participantes do Exame Nacional do Ensino Médio (ENEM), permitindo analisar aspectos como o desempenho dos alunos nas diferentes áreas da prova e suas características socioeconômicas.

A partir desses dados, busca-se responder à seguinte questão: A situação socioeconômica dos alunos influencia em suas notas?

O conjunto de dados disponibilizado foi previamente extraído e posteriormente organizado, mantendo apenas as colunas consideradas relevantes para o desenvolvimento do estudo e para a análise da relação entre as condições socioeconômicas e o desempenho dos participantes.

\section{Dicionário de Variáveis}

O dicionário de variáveis apresenta as principais colunas utilizadas no conjunto de dados, acompanhadas de uma breve descrição de seu significado, seu respectivo tipo e seus valores ou unidades.

Para facilitar a organização e a apresentação visual da tabela, alguns valores foram resumidos mas pelo menos o valor mínimo e máximo presente neles é explicado na linha específica da tabela como por exemplo na Q07 que tem 17 respostas possíveis da renda mensal da família e portanto foi resumida ao valor inicial de nenhuma renda e acima de Acima de 30.360,00 reais. Entretanto, no artigo final, as categorias e os valores completos serão apresentados e devidamente detalhados.


\begin{table}[ht]
\centering
 
\caption*{\textbf{Dicionário de dados}}
 
\vspace{0.2cm}
 
\renewcommand{\arraystretch}{1.8} %mais linhas
 
\begin{tabularx}{\textwidth}{
>{\centering\arraybackslash}p{3cm} % tamanho Variável
>{\raggedright\arraybackslash}X     % Tamanho Descrição
>{\centering\arraybackslash}p{4cm}  % Tamanho Tipo
>{\centering\arraybackslash}p{4cm}  % Tamanho Valores
}
\toprule
\textbf{Variavel} &
\textbf{Descricao} &
\textbf{Tipo} &
\textbf{Valores / unidade} \\
\midrule
 
 
NU\_INSCRICAO&Número da inscrição do aluno&Qualitativa nominal& Número de 12 dígitos \\[0.3cm]
NU\_ANO&Ano de realização da prova&Quantitativa discreta& Anos de 2021 à 2025 \\[0.3cm]
TP\_COR\_RACA& Identificação de raça& Qualitativa nominal & Categorias de Raça/Cor & \\[0.3cm]
NU\_NOTA\_CN& Nota de ciências naturais & Quantitativa contínua & 0 a 1000 pontos \\[0.3cm]
NU\_NOTA\_CH& Nota de ciências humanas& Quantitativa contínua & 0 a 1000 pontos \\[0.3cm]
NU\_NOTA\_LC& Nota de linguagens e codigos & Quantitativa contínua & 0 a 1000 pontos \\[0.3cm]
NU\_NOTA\_MT& Nota de matemática &Quantitativa contínua & 0 a 1000 pontos \\[0.3cm]
NU\_NOTA\_RED& Nota da redação & Quantitativa contínua & 0 a 1000 pontos \\[0.3cm]
Q006& O aluno possui renda? & Quantitativa nominal & A = Não; B = Sim  \\[0.3cm]
Q007& Qual a renda mensal (Reais) da família do aluno? & Qualitativa ordinal & Valores de A = Nenhuma até Q = Acima de 30.360,00 \\[0.3cm]
Q020& Na casa do aluno, existe acesso à internet por rede wi-fi? & Qualitativa nominal & A = Não; B = Sim \\[0.3cm]
Q021& Na casa do aluno, existe computador/notebook? & Qualitativa nominal & Categorias de A = Não até E = Sim, quatro ou mais \\[0.3cm]

\bottomrule
\end{tabularx}
 
\end{table}
 
\end{document}
